% !TeX program = xelatex
% 中英文块级字段审稿回复信模板
% 推荐连续编译两次：xelatex Review_Response_Template.tex

\documentclass[10pt,a4paper]{article}

\usepackage[
  a4paper,
  top=1in,
  bottom=1in,
  left=1.25in,
  right=1.25in,
  headheight=0pt,
  headsep=0pt,
  footskip=0pt
]{geometry}
\usepackage{fontspec}
\usepackage{xeCJK}
\setmainfont{Times New Roman}
\setCJKmainfont{SimSun}[AutoFakeBold=2]
\usepackage{amsmath,amssymb}
\usepackage{graphicx}
\usepackage{booktabs,tabularx,array}
\usepackage{etoolbox}
\usepackage{needspace}
\usepackage{caption}
\usepackage[hidelinks]{hyperref}

\pagestyle{empty}
\setlength{\parindent}{0pt}
\setlength{\parskip}{0pt}
\setlength{\emergencystretch}{2em}
\linespread{1.0}\selectfont
\sloppy
\clubpenalty=10000
\widowpenalty=10000
\displaywidowpenalty=10000
\hyphenpenalty=10000
\exhyphenpenalty=10000

\captionsetup{
  font=normalsize,
  labelfont=normalfont,
  labelsep=period,
  justification=centering,
  singlelinecheck=true,
  skip=5pt
}

% =============================================================================
% 只需修改这一行即可切换显示模式：
%   en   = 仅英文
%   cn   = 仅中文
%   both = 英文整块 + 对应中文整块
% 命令行临时切换示例：
% xelatex "\def\ResponseMode{cn}% !TeX program = xelatex
% 中英文块级字段审稿回复信模板
% 推荐连续编译两次：xelatex Review_Response_Template.tex

\documentclass[10pt,a4paper]{article}

\usepackage[
  a4paper,
  top=1in,
  bottom=1in,
  left=1.25in,
  right=1.25in,
  headheight=0pt,
  headsep=0pt,
  footskip=0pt
]{geometry}
\usepackage{fontspec}
\usepackage{xeCJK}
\setmainfont{Times New Roman}
\setCJKmainfont{SimSun}[AutoFakeBold=2]
\usepackage{amsmath,amssymb}
\usepackage{graphicx}
\usepackage{booktabs,tabularx,array}
\usepackage{etoolbox}
\usepackage{needspace}
\usepackage{caption}
\usepackage[hidelinks]{hyperref}

\pagestyle{empty}
\setlength{\parindent}{0pt}
\setlength{\parskip}{0pt}
\setlength{\emergencystretch}{2em}
\linespread{1.0}\selectfont
\sloppy
\clubpenalty=10000
\widowpenalty=10000
\displaywidowpenalty=10000
\hyphenpenalty=10000
\exhyphenpenalty=10000

\captionsetup{
  font=normalsize,
  labelfont=normalfont,
  labelsep=period,
  justification=centering,
  singlelinecheck=true,
  skip=5pt
}

% =============================================================================
% 只需修改这一行即可切换显示模式：
%   en   = 仅英文
%   cn   = 仅中文
%   both = 英文整块 + 对应中文整块
% 命令行临时切换示例：
% xelatex "\def\ResponseMode{cn}% !TeX program = xelatex
% 中英文块级字段审稿回复信模板
% 推荐连续编译两次：xelatex Review_Response_Template.tex

\documentclass[10pt,a4paper]{article}

\usepackage[
  a4paper,
  top=1in,
  bottom=1in,
  left=1.25in,
  right=1.25in,
  headheight=0pt,
  headsep=0pt,
  footskip=0pt
]{geometry}
\usepackage{fontspec}
\usepackage{xeCJK}
\setmainfont{Times New Roman}
\setCJKmainfont{SimSun}[AutoFakeBold=2]
\usepackage{amsmath,amssymb}
\usepackage{graphicx}
\usepackage{booktabs,tabularx,array}
\usepackage{etoolbox}
\usepackage{needspace}
\usepackage{caption}
\usepackage[hidelinks]{hyperref}

\pagestyle{empty}
\setlength{\parindent}{0pt}
\setlength{\parskip}{0pt}
\setlength{\emergencystretch}{2em}
\linespread{1.0}\selectfont
\sloppy
\clubpenalty=10000
\widowpenalty=10000
\displaywidowpenalty=10000
\hyphenpenalty=10000
\exhyphenpenalty=10000

\captionsetup{
  font=normalsize,
  labelfont=normalfont,
  labelsep=period,
  justification=centering,
  singlelinecheck=true,
  skip=5pt
}

% =============================================================================
% 只需修改这一行即可切换显示模式：
%   en   = 仅英文
%   cn   = 仅中文
%   both = 英文整块 + 对应中文整块
% 命令行临时切换示例：
% xelatex "\def\ResponseMode{cn}% !TeX program = xelatex
% 中英文块级字段审稿回复信模板
% 推荐连续编译两次：xelatex Review_Response_Template.tex

\documentclass[10pt,a4paper]{article}

\usepackage[
  a4paper,
  top=1in,
  bottom=1in,
  left=1.25in,
  right=1.25in,
  headheight=0pt,
  headsep=0pt,
  footskip=0pt
]{geometry}
\usepackage{fontspec}
\usepackage{xeCJK}
\setmainfont{Times New Roman}
\setCJKmainfont{SimSun}[AutoFakeBold=2]
\usepackage{amsmath,amssymb}
\usepackage{graphicx}
\usepackage{booktabs,tabularx,array}
\usepackage{etoolbox}
\usepackage{needspace}
\usepackage{caption}
\usepackage[hidelinks]{hyperref}

\pagestyle{empty}
\setlength{\parindent}{0pt}
\setlength{\parskip}{0pt}
\setlength{\emergencystretch}{2em}
\linespread{1.0}\selectfont
\sloppy
\clubpenalty=10000
\widowpenalty=10000
\displaywidowpenalty=10000
\hyphenpenalty=10000
\exhyphenpenalty=10000

\captionsetup{
  font=normalsize,
  labelfont=normalfont,
  labelsep=period,
  justification=centering,
  singlelinecheck=true,
  skip=5pt
}

% =============================================================================
% 只需修改这一行即可切换显示模式：
%   en   = 仅英文
%   cn   = 仅中文
%   both = 英文整块 + 对应中文整块
% 命令行临时切换示例：
% xelatex "\def\ResponseMode{cn}\input{Review_Response_Template.tex}"
% =============================================================================
\providecommand{\ResponseMode}{both}

\newif\ifShowEnglish
\newif\ifShowChinese

\newcommand{\ResponseDisplayMode}[1]{%
  \ShowEnglishfalse
  \ShowChinesefalse
  \ifstrequal{#1}{en}{%
    \ShowEnglishtrue
  }{%
    \ifstrequal{#1}{cn}{%
      \ShowChinesetrue
    }{%
      \ifstrequal{#1}{both}{%
        \ShowEnglishtrue
        \ShowChinesetrue
      }{%
        \PackageError{ReviewResponse}{Unknown display mode '#1'}%
          {Use en, cn, or both.}%
      }%
    }%
  }%
}
\expandafter\ResponseDisplayMode\expandafter{\ResponseMode}

\newcommand{\ModeInline}[2]{%
  \ifShowEnglish #1\ifShowChinese\ /\ \fi\fi
  \ifShowChinese #2\fi
}

\renewcommand{\figurename}{\ModeInline{Figure}{图}}
\renewcommand{\tablename}{\ModeInline{Table}{表}}

% =============================================================================
%                         用户编辑区 1：信件整体字段
% 只需在这里填写一次，标题、开场说明和署名会自动同步替换。
% 未确定的内容暂时保留为普通黑色 xxx。
% =============================================================================

% 论文英文标题
\newcommand{\MetaTitleEN}{xxx}
% 论文中文标题
\newcommand{\MetaTitleCN}{xxx}
% 稿件编号（中英文共用，只填一次）
\newcommand{\MetaManuscriptID}{xxx}
% 通讯作者英文姓名
\newcommand{\MetaCorrespondingAuthorEN}{xxx}
% 通讯作者中文姓名
\newcommand{\MetaCorrespondingAuthorCN}{xxx}

\newcommand{\FieldSalutationEN}{Dear Editor and Reviewers,}
\newcommand{\FieldSalutationCN}{尊敬的编辑和审稿专家：}

% 每条回复的固定首句，由模板自动添加，无需在 ResponseEN / ResponseCN 中重复填写。
\newcommand{\ResponseOpeningEN}{We thank the reviewer for these comments.}
\newcommand{\ResponseOpeningCN}{感谢审稿人提出这些意见。}

\newcommand{\FieldOpeningEN}{%
  Thank you for your careful evaluation and constructive comments. We have
  carefully considered each comment and revised the manuscript accordingly. Our
  point-by-point responses are provided below. The reviewers' comments are
  reproduced in bold, followed by our responses. All changes in the revised
  manuscript are highlighted in red.%
}

\newcommand{\FieldOpeningCN}{%
  感谢您对本稿件的认真评阅和建设性意见。我们已逐条认真研究并据此修订稿件。下文
  给出逐点回复。审稿意见以粗体显示，随后给出我们的回复。修订稿中的所有改动均以
  红色标出。%
}

\newcommand{\FieldSignOffEN}{%
  With kindest regards,\par
  Yours Sincerely,\par
  \MetaCorrespondingAuthorEN%
}

\newcommand{\FieldSignOffCN}{%
  此致\par
  敬礼！\par
  \MetaCorrespondingAuthorCN%
}

% =============================================================================
%                         以下为排版逻辑，通常无需修改
% =============================================================================

\newcommand{\IndentedParagraph}[1]{%
  \par\begingroup
  \setlength{\parindent}{21pt}%
  \indent #1\par
  \endgroup
}

\newcommand{\MakeLetterHeader}{%
  \begin{center}
    \bfseries\fontsize{12pt}{14.4pt}\selectfont
    \ifShowEnglish Revisions to the paper titled \MetaTitleEN\fi
    \ifShowEnglish\ifShowChinese\\[0.35em]\fi\fi
    \ifShowChinese 关于论文《\MetaTitleCN》的修改说明\fi
  \end{center}
  \vspace{1.2\baselineskip}

  \ifShowEnglish\noindent Manuscript Number: \MetaManuscriptID\par\fi
  \ifShowChinese\noindent 稿件编号：\MetaManuscriptID\par\fi
  \vspace{1.9\baselineskip}

  \ifShowEnglish \FieldSalutationEN\par\fi
  \ifShowEnglish\ifShowChinese\smallskip\fi\fi
  \ifShowChinese \FieldSalutationCN\par\fi
  \medskip

  \ifShowEnglish\IndentedParagraph{\FieldOpeningEN}\fi
  \ifShowEnglish\ifShowChinese\vspace{0.45\baselineskip}\fi\fi
  \ifShowChinese\IndentedParagraph{\FieldOpeningCN}\fi

  \vspace{0.65\baselineskip}
  \ifShowEnglish \FieldSignOffEN\par\fi
  \ifShowEnglish\ifShowChinese\vspace{0.45\baselineskip}\fi\fi
  \ifShowChinese \FieldSignOffCN\par\fi
}

\newcounter{reviewcomment}

\newcommand{\EditorSection}{%
  \vspace{1.35\baselineskip}
  \noindent\textbf{\ModeInline{Response to the Editor:}{对编辑意见的回复：}}\par
  \vspace{0.75\baselineskip}
  \setcounter{reviewcomment}{0}%
}

\newcommand{\ReviewerSection}[1]{%
  \vspace{1.35\baselineskip}
  \noindent\textbf{\ModeInline{Response to Reviewer \##1:}{对审稿人 \##1 的回复：}}\par
  \vspace{0.75\baselineskip}
  \setcounter{reviewcomment}{0}%
}

% 审稿意见：先显示完整英文块，再显示完整中文块；不做逐句交替。
\NewDocumentCommand{\RenderCommentBlock}{m +m +m}{%
  \refstepcounter{reviewcomment}%
  \Needspace{7\baselineskip}%
  \par\vspace{0.9\baselineskip}
  \noindent\hypertarget{#1}{}%
  \begingroup\bfseries
  \arabic{reviewcomment}.\ %
  \ifShowEnglish #2\fi
  \ifShowEnglish\ifShowChinese\par\smallskip\noindent\hspace*{1.35em}\fi\fi
  \ifShowChinese #3\fi
  \par\endgroup
  \nobreak\smallskip
}

% 作者回复：先显示完整英文块，再显示完整中文块；不做逐句交替。
\NewDocumentCommand{\RenderResponseBlock}{+m +m}{%
  \Needspace{4\baselineskip}%
  \noindent\textbf{\ModeInline{Response:}{回复：}}\par
  \nobreak
  \ifShowEnglish
    \noindent\mbox{\textgreater\kern0.08em\textgreater}\
    \ResponseOpeningEN\space #1\par
  \fi
  \ifShowEnglish\ifShowChinese\smallskip\fi\fi
  \ifShowChinese
    \noindent\mbox{\textgreater\kern0.08em\textgreater}\
    \ResponseOpeningCN\space #2\par
  \fi
}

\NewDocumentCommand{\RenderRevisedText}{+m +m}{%
  \ifShowEnglish
    \par\begingroup
    \setlength{\leftskip}{21pt}\setlength{\rightskip}{21pt}%
    \textit{Revised text: }``#1''\par
    \endgroup
  \fi
  \ifShowEnglish\ifShowChinese\smallskip\fi\fi
  \ifShowChinese
    \par\begingroup
    \setlength{\leftskip}{21pt}\setlength{\rightskip}{21pt}%
    修订文本：“#2”\par
    \endgroup
  \fi
}

\NewDocumentCommand{\RenderLocation}{+m +m}{%
  \ifShowEnglish\par\noindent The corresponding change appears in #1.\par\fi
  \ifShowChinese\par\noindent 相应修改位于#2。\par\fi
}

% 自动位置字段：重复使用 \At{页码}{行号或行号范围} 即可添加多个位置。
% 同一页出现多个 \At 时会自动合并；多个页面会自动补充 and / 以及。
\ExplSyntaxOn
\seq_new:N  \l__rr_location_pages_seq
\seq_new:N  \l__rr_location_lines_seq
\prop_new:N \l__rr_location_lines_prop
\bool_new:N \l__rr_location_singular_line_bool
\int_new:N  \l__rr_location_page_count_int
\int_new:N  \l__rr_location_line_count_int

\cs_generate_variant:Nn \seq_set_split:Nnn { NnV }

\cs_new_protected:Npn \__rr_location_collect:nn #1#2
  {
    \prop_get:NnNTF \l__rr_location_lines_prop {#1} \l_tmpa_tl
      {
        \prop_put:Nnx \l__rr_location_lines_prop {#1}
          { \exp_not:V \l_tmpa_tl ; \exp_not:n {#2} }
      }
      {
        \seq_put_right:Nn \l__rr_location_pages_seq {#1}
        \prop_put:Nnn \l__rr_location_lines_prop {#1} {#2}
      }
  }

\cs_new_protected:Npn \__rr_location_prepare:n #1
  {
    \seq_clear:N \l__rr_location_pages_seq
    \prop_clear:N \l__rr_location_lines_prop
    \cs_set_protected:Npn \At ##1##2
      { \__rr_location_collect:nn {##1} {##2} }
    #1
  }

\cs_new_protected:Npn \__rr_location_set_lines_for_page:n #1
  {
    \prop_get:NnN \l__rr_location_lines_prop {#1} \l_tmpa_tl
    \seq_set_split:NnV \l__rr_location_lines_seq {;} \l_tmpa_tl
    \seq_set_map:NNn \l__rr_location_lines_seq \l__rr_location_lines_seq
      { \tl_trim_spaces:n {##1} }
  }

\cs_new_protected:Npn \__rr_location_set_line_number:n #1
  {
    \__rr_location_set_lines_for_page:n {#1}
    \bool_set_false:N \l__rr_location_singular_line_bool
    \int_compare:nNnT { \seq_count:N \l__rr_location_lines_seq } = {1}
      {
        \tl_set:Nx \l_tmpa_tl
          { \seq_item:Nn \l__rr_location_lines_seq {1} }
        \tl_if_in:VnF \l_tmpa_tl {--}
          {
            \tl_if_in:VnF \l_tmpa_tl {,}
              { \bool_set_true:N \l__rr_location_singular_line_bool }
          }
      }
  }

\cs_new_protected:Npn \__rr_location_en_line_item:n #1
  {
    \tl_set:Nn \l_tmpa_tl {#1}
    \regex_replace_all:nnN { \s* , \s* } { ~and~ } \l_tmpa_tl
    \tl_use:N \l_tmpa_tl
  }

\cs_new_protected:Npn \__rr_location_en_lines:
  {
    \int_set:Nn \l__rr_location_line_count_int
      { \seq_count:N \l__rr_location_lines_seq }
    \seq_map_indexed_inline:Nn \l__rr_location_lines_seq
      {
        \int_compare:nNnT {##1} > {1}
          {
            \int_compare:nNnTF {\l__rr_location_line_count_int} = {2}
              { ~and~ }
              {
                \int_compare:nNnTF {##1} = {\l__rr_location_line_count_int}
                  { ,~and~ }
                  { ,~ }
              }
          }
        \__rr_location_en_line_item:n {##2}
      }
  }

\cs_new_protected:Npn \__rr_location_en_page:n #1
  {
    \__rr_location_set_line_number:n {#1}
    \tl_if_in:nnTF {#1} {--} {pages} {page}\space #1
    \int_compare:nNnTF {\l__rr_location_page_count_int} = {1}
      {
        ,\space
        \bool_if:NTF \l__rr_location_singular_line_bool {line} {lines}
        \space\__rr_location_en_lines:
      }
      {
        \space(
        \bool_if:NTF \l__rr_location_singular_line_bool {line} {lines}
        \space\__rr_location_en_lines:)
      }
  }

\cs_new_protected:Npn \__rr_location_en_pages:
  {
    \int_set:Nn \l__rr_location_page_count_int
      { \seq_count:N \l__rr_location_pages_seq }
    \seq_map_indexed_inline:Nn \l__rr_location_pages_seq
      {
        \int_compare:nNnT {##1} > {1}
          {
            \int_compare:nNnTF {\l__rr_location_page_count_int} = {2}
              { \space and\space }
              {
                \int_compare:nNnTF {##1} = {\l__rr_location_page_count_int}
                  { ,\space and\space }
                  { ,\space }
              }
          }
        \__rr_location_en_page:n {##2}
      }
  }

\cs_new_protected:Npn \__rr_location_cn_lines:
  {
    \int_set:Nn \l__rr_location_line_count_int
      { \seq_count:N \l__rr_location_lines_seq }
    \seq_map_indexed_inline:Nn \l__rr_location_lines_seq
      {
        \int_compare:nNnT {##1} > {1}
          {
            \int_compare:nNnTF {##1} = {\l__rr_location_line_count_int}
              { 和 }
              { 、 }
          }
        第~##2~行
      }
  }

\cs_new_protected:Npn \__rr_location_cn_page:n #1
  {
    \__rr_location_set_lines_for_page:n {#1}
    第~#1~页
    \int_compare:nNnTF {\l__rr_location_page_count_int} = {1}
      { \__rr_location_cn_lines: }
      { （\__rr_location_cn_lines:） }
  }

\cs_new_protected:Npn \__rr_location_cn_pages:
  {
    \int_set:Nn \l__rr_location_page_count_int
      { \seq_count:N \l__rr_location_pages_seq }
    \seq_map_indexed_inline:Nn \l__rr_location_pages_seq
      {
        \int_compare:nNnT {##1} > {1}
          {
            \int_compare:nNnTF {\l__rr_location_page_count_int} = {2}
              { 以及 }
              {
                \int_compare:nNnTF {##1} = {\l__rr_location_page_count_int}
                  { 以及 }
                  { 、 }
              }
          }
        \__rr_location_cn_page:n {##2}
      }
  }

\cs_new_protected:Npn \__rr_render_auto_location:n #1
  {
    \__rr_location_prepare:n {#1}
    \ifShowEnglish
      \par\noindent
      \int_compare:nNnTF { \seq_count:N \l__rr_location_pages_seq } = {1}
        { The~corresponding~change~appears~in~the~revised~manuscript~on~ }
        { The~corresponding~changes~appear~in~the~revised~manuscript~on~ }
      \__rr_location_en_pages:.\par
    \fi
    \ifShowChinese
      \par\noindent
      相应修改位于修订稿
      \__rr_location_cn_pages:。\par
    \fi
  }
\NewDocumentCommand{\RenderAutoLocation}{+m}{%
  \__rr_render_auto_location:n {#1}%
}
\ExplSyntaxOff

% -----------------------------------------------------------------------------
% 一条审稿意见的友好字段环境
%
% \begin{ReviewItem}{稳定编号}
%   \CommentEN{完整英文意见}
%   \CommentCN{完整中文意见}
%   \ResponseEN{英文固定首句之后的回复正文}
%   \ResponseCN{中文固定首句之后的回复正文}
%   \RevisedEN{可选：英文修订原文}
%   \RevisedCN{可选：中文对照}
%   \Evidence{可选：图、表或公式}
%   \Location{\At{8}{210--225}\At{8}{244--260}\At{10}{301}\At{12}{365}}
%   特殊位置无法用页码和行号表达时，改用 \LocationEN 和 \LocationCN 手写。
% \end{ReviewItem}
% -----------------------------------------------------------------------------
\NewDocumentEnvironment{ReviewItem}{m +b}{%
  \begingroup
  \long\def\RRCommentEN{xxx}%
  \long\def\RRCommentCN{xxx}%
  \long\def\RRResponseEN{xxx}%
  \long\def\RRResponseCN{xxx}%
  \long\def\RRRevisedEN{}%
  \long\def\RRRevisedCN{}%
  \long\def\RREvidence{}%
  \long\def\RRAutoLocation{}%
  \long\def\RRLocationEN{xxx}%
  \long\def\RRLocationCN{xxx}%

  \long\def\CommentEN##1{\long\def\RRCommentEN{##1}}%
  \long\def\CommentCN##1{\long\def\RRCommentCN{##1}}%
  \long\def\ResponseEN##1{\long\def\RRResponseEN{##1}}%
  \long\def\ResponseCN##1{\long\def\RRResponseCN{##1}}%
  \long\def\RevisedEN##1{\long\def\RRRevisedEN{##1}}%
  \long\def\RevisedCN##1{\long\def\RRRevisedCN{##1}}%
  \long\def\Evidence##1{\long\def\RREvidence{##1}}%
  \long\def\Location##1{\long\def\RRAutoLocation{##1}}%
  \long\def\LocationEN##1{\long\def\RRLocationEN{##1}}%
  \long\def\LocationCN##1{\long\def\RRLocationCN{##1}}%

  #2%

  \RenderCommentBlock{#1}{\RRCommentEN}{\RRCommentCN}%
  \RenderResponseBlock{\RRResponseEN}{\RRResponseCN}%
  \ifdefempty{\RRRevisedEN}{%
    \ifdefempty{\RRRevisedCN}{}{\RenderRevisedText{\RRRevisedEN}{\RRRevisedCN}}%
  }{\RenderRevisedText{\RRRevisedEN}{\RRRevisedCN}}%
  \ifdefempty{\RREvidence}{}{\RREvidence}%
  \ifdefempty{\RRAutoLocation}{%
    \RenderLocation{\RRLocationEN}{\RRLocationCN}%
  }{\RenderAutoLocation{\RRAutoLocation}}%
  \endgroup
}{}

\NewDocumentCommand{\FigurePlaceholder}{O{0.66\linewidth} m m}{%
  \begin{figure}[htbp]
    \centering
    \fbox{\rule{0pt}{30mm}\rule{#1}{0pt}}
    \caption{\ModeInline{#2}{#3}}
  \end{figure}
}

\newcolumntype{Y}{>{\raggedright\arraybackslash}X}

\begin{document}

\MakeLetterHeader

% =============================================================================
%                      用户编辑区 2：逐条审稿意见与回复
% 每个字段均填写一整段或多段内容，不要逐句拆分。
% =============================================================================

\ReviewerSection{1}

\begin{ReviewItem}{R1.1}
  % 审稿意见英文原文（完整粘贴，不要改写）
  \CommentEN{xxx}
  % 审稿意见中文译文（完整填写）
  \CommentCN{xxx}
  % 英文回复正文；固定首句由模板自动添加，只需填写首句之后的内容
  \ResponseEN{%
    Specifically, xxx. Accordingly, we have revised xxx to improve xxx.%
  }
  % 中文回复正文；固定首句由模板自动添加，只需填写首句之后的内容
  \ResponseCN{%
    具体而言，xxx。因此，我们已对 xxx 进行了修改，以进一步完善 xxx。%
  }
  % 可选：修订稿中的英文原文；不需要引用时可删除本行
  \RevisedEN{xxx}
  % 可选：修订原文的中文对照；不需要引用时可删除本行
  \RevisedCN{xxx}
  % 修改位置：每个 \At{页码}{行号} 表示一处；重复 \At 即可添加多个位置
  % 单个位置示例：\Location{\At{8}{210--225}}
  % 多个位置示例（相同页码会自动合并）：
  % \Location{\At{8}{210--225}\At{8}{244--260}\At{10}{301}\At{12}{365}}
  \Location{\At{xxx}{xxx}}
\end{ReviewItem}

\begin{ReviewItem}{R1.2}
  % 审稿意见英文原文（完整粘贴，不要改写）
  \CommentEN{xxx}
  % 审稿意见中文译文（完整填写）
  \CommentCN{xxx}
  % 英文回复正文；固定首句由模板自动添加
  \ResponseEN{%
    We agree that the original presentation did not make this point sufficiently
    clear. We have therefore clarified xxx and added xxx to support this revision.%
  }
  % 中文回复正文；固定首句由模板自动添加
  \ResponseCN{%
    我们同意原稿对此说明得不够清楚。因此，我们已进一步澄清 xxx，并补充 xxx 作为
    支撑。%
  }
  % 可选：回复中需要展示的图、表或公式；不需要时可删除整个 Evidence 字段
  \Evidence{%
    % 替换为真实图片时使用：
    % \includegraphics[width=0.66\linewidth]{figures/your-figure.pdf}
    \FigurePlaceholder
      {xxx}
      {xxx}%
  }
  % 修改位置：重复 \At{页码}{行号或行号范围} 可添加多个位置
  \Location{\At{xxx}{xxx}}
\end{ReviewItem}

\begin{ReviewItem}{R1.3}
  % 审稿意见英文原文（完整粘贴，不要改写）
  \CommentEN{xxx}
  % 审稿意见中文译文（完整填写）
  \CommentCN{xxx}
  % 英文回复正文；固定首句由模板自动添加
  \ResponseEN{%
    The requested comparison has been added and is summarized in
    Table~\ref{tab:comparison}.%
  }
  % 中文回复正文；固定首句由模板自动添加
  \ResponseCN{%
    我们已补充所要求的比较，结果汇总于表~\ref{tab:comparison}。%
  }
  % 可选：定量比较表；将表题、方法名、指标名和数值中的 xxx 替换为实际内容
  \Evidence{%
    \begin{table}[htbp]
      \centering
      \caption{\ModeInline
        {xxx}
        {xxx}}
      \label{tab:comparison}
      \begin{tabularx}{0.86\linewidth}{@{}Y>{\centering\arraybackslash}p{2.6cm}>{\centering\arraybackslash}p{2.6cm}@{}}
        \toprule
        \ModeInline{Method}{方法} & \ModeInline{Metric 1}{指标 1} & \ModeInline{Metric 2}{指标 2} \\
        \midrule
        \ModeInline{Proposed method}{本文方法} & xxx & xxx \\
        \ModeInline{Baseline method}{基线方法} & xxx & xxx \\
        \bottomrule
      \end{tabularx}
    \end{table}%
  }
  % 修改位置：重复 \At{页码}{行号或行号范围} 可添加多个位置
  \Location{\At{xxx}{xxx}}
\end{ReviewItem}

\ReviewerSection{2}

\begin{ReviewItem}{R2.1}
  % 审稿意见英文原文（完整粘贴，不要改写）
  \CommentEN{xxx}
  % 审稿意见中文译文（完整填写）
  \CommentCN{xxx}
  % 英文回复正文；固定首句由模板自动添加
  \ResponseEN{%
    In response, we have revised xxx and added xxx to substantiate this point.%
  }
  % 中文回复正文；固定首句由模板自动添加
  \ResponseCN{%
    作为回应，我们已修改 xxx，并补充 xxx 以提供支撑。%
  }
  % 修改位置：重复 \At{页码}{行号或行号范围} 可添加多个位置
  \Location{\At{xxx}{xxx}}
\end{ReviewItem}

% 添加新意见：复制一个完整的 ReviewItem 环境。
% 添加新审稿人：复制 ReviewerSection，并从新的 ReviewItem 开始。

\end{document}
"
% =============================================================================
\providecommand{\ResponseMode}{both}

\newif\ifShowEnglish
\newif\ifShowChinese

\newcommand{\ResponseDisplayMode}[1]{%
  \ShowEnglishfalse
  \ShowChinesefalse
  \ifstrequal{#1}{en}{%
    \ShowEnglishtrue
  }{%
    \ifstrequal{#1}{cn}{%
      \ShowChinesetrue
    }{%
      \ifstrequal{#1}{both}{%
        \ShowEnglishtrue
        \ShowChinesetrue
      }{%
        \PackageError{ReviewResponse}{Unknown display mode '#1'}%
          {Use en, cn, or both.}%
      }%
    }%
  }%
}
\expandafter\ResponseDisplayMode\expandafter{\ResponseMode}

\newcommand{\ModeInline}[2]{%
  \ifShowEnglish #1\ifShowChinese\ /\ \fi\fi
  \ifShowChinese #2\fi
}

\renewcommand{\figurename}{\ModeInline{Figure}{图}}
\renewcommand{\tablename}{\ModeInline{Table}{表}}

% =============================================================================
%                         用户编辑区 1：信件整体字段
% 只需在这里填写一次，标题、开场说明和署名会自动同步替换。
% 未确定的内容暂时保留为普通黑色 xxx。
% =============================================================================

% 论文英文标题
\newcommand{\MetaTitleEN}{xxx}
% 论文中文标题
\newcommand{\MetaTitleCN}{xxx}
% 稿件编号（中英文共用，只填一次）
\newcommand{\MetaManuscriptID}{xxx}
% 通讯作者英文姓名
\newcommand{\MetaCorrespondingAuthorEN}{xxx}
% 通讯作者中文姓名
\newcommand{\MetaCorrespondingAuthorCN}{xxx}

\newcommand{\FieldSalutationEN}{Dear Editor and Reviewers,}
\newcommand{\FieldSalutationCN}{尊敬的编辑和审稿专家：}

% 每条回复的固定首句，由模板自动添加，无需在 ResponseEN / ResponseCN 中重复填写。
\newcommand{\ResponseOpeningEN}{We thank the reviewer for these comments.}
\newcommand{\ResponseOpeningCN}{感谢审稿人提出这些意见。}

\newcommand{\FieldOpeningEN}{%
  Thank you for your careful evaluation and constructive comments. We have
  carefully considered each comment and revised the manuscript accordingly. Our
  point-by-point responses are provided below. The reviewers' comments are
  reproduced in bold, followed by our responses. All changes in the revised
  manuscript are highlighted in red.%
}

\newcommand{\FieldOpeningCN}{%
  感谢您对本稿件的认真评阅和建设性意见。我们已逐条认真研究并据此修订稿件。下文
  给出逐点回复。审稿意见以粗体显示，随后给出我们的回复。修订稿中的所有改动均以
  红色标出。%
}

\newcommand{\FieldSignOffEN}{%
  With kindest regards,\par
  Yours Sincerely,\par
  \MetaCorrespondingAuthorEN%
}

\newcommand{\FieldSignOffCN}{%
  此致\par
  敬礼！\par
  \MetaCorrespondingAuthorCN%
}

% =============================================================================
%                         以下为排版逻辑，通常无需修改
% =============================================================================

\newcommand{\IndentedParagraph}[1]{%
  \par\begingroup
  \setlength{\parindent}{21pt}%
  \indent #1\par
  \endgroup
}

\newcommand{\MakeLetterHeader}{%
  \begin{center}
    \bfseries\fontsize{12pt}{14.4pt}\selectfont
    \ifShowEnglish Revisions to the paper titled \MetaTitleEN\fi
    \ifShowEnglish\ifShowChinese\\[0.35em]\fi\fi
    \ifShowChinese 关于论文《\MetaTitleCN》的修改说明\fi
  \end{center}
  \vspace{1.2\baselineskip}

  \ifShowEnglish\noindent Manuscript Number: \MetaManuscriptID\par\fi
  \ifShowChinese\noindent 稿件编号：\MetaManuscriptID\par\fi
  \vspace{1.9\baselineskip}

  \ifShowEnglish \FieldSalutationEN\par\fi
  \ifShowEnglish\ifShowChinese\smallskip\fi\fi
  \ifShowChinese \FieldSalutationCN\par\fi
  \medskip

  \ifShowEnglish\IndentedParagraph{\FieldOpeningEN}\fi
  \ifShowEnglish\ifShowChinese\vspace{0.45\baselineskip}\fi\fi
  \ifShowChinese\IndentedParagraph{\FieldOpeningCN}\fi

  \vspace{0.65\baselineskip}
  \ifShowEnglish \FieldSignOffEN\par\fi
  \ifShowEnglish\ifShowChinese\vspace{0.45\baselineskip}\fi\fi
  \ifShowChinese \FieldSignOffCN\par\fi
}

\newcounter{reviewcomment}

\newcommand{\EditorSection}{%
  \vspace{1.35\baselineskip}
  \noindent\textbf{\ModeInline{Response to the Editor:}{对编辑意见的回复：}}\par
  \vspace{0.75\baselineskip}
  \setcounter{reviewcomment}{0}%
}

\newcommand{\ReviewerSection}[1]{%
  \vspace{1.35\baselineskip}
  \noindent\textbf{\ModeInline{Response to Reviewer \##1:}{对审稿人 \##1 的回复：}}\par
  \vspace{0.75\baselineskip}
  \setcounter{reviewcomment}{0}%
}

% 审稿意见：先显示完整英文块，再显示完整中文块；不做逐句交替。
\NewDocumentCommand{\RenderCommentBlock}{m +m +m}{%
  \refstepcounter{reviewcomment}%
  \Needspace{7\baselineskip}%
  \par\vspace{0.9\baselineskip}
  \noindent\hypertarget{#1}{}%
  \begingroup\bfseries
  \arabic{reviewcomment}.\ %
  \ifShowEnglish #2\fi
  \ifShowEnglish\ifShowChinese\par\smallskip\noindent\hspace*{1.35em}\fi\fi
  \ifShowChinese #3\fi
  \par\endgroup
  \nobreak\smallskip
}

% 作者回复：先显示完整英文块，再显示完整中文块；不做逐句交替。
\NewDocumentCommand{\RenderResponseBlock}{+m +m}{%
  \Needspace{4\baselineskip}%
  \noindent\textbf{\ModeInline{Response:}{回复：}}\par
  \nobreak
  \ifShowEnglish
    \noindent\mbox{\textgreater\kern0.08em\textgreater}\
    \ResponseOpeningEN\space #1\par
  \fi
  \ifShowEnglish\ifShowChinese\smallskip\fi\fi
  \ifShowChinese
    \noindent\mbox{\textgreater\kern0.08em\textgreater}\
    \ResponseOpeningCN\space #2\par
  \fi
}

\NewDocumentCommand{\RenderRevisedText}{+m +m}{%
  \ifShowEnglish
    \par\begingroup
    \setlength{\leftskip}{21pt}\setlength{\rightskip}{21pt}%
    \textit{Revised text: }``#1''\par
    \endgroup
  \fi
  \ifShowEnglish\ifShowChinese\smallskip\fi\fi
  \ifShowChinese
    \par\begingroup
    \setlength{\leftskip}{21pt}\setlength{\rightskip}{21pt}%
    修订文本：“#2”\par
    \endgroup
  \fi
}

\NewDocumentCommand{\RenderLocation}{+m +m}{%
  \ifShowEnglish\par\noindent The corresponding change appears in #1.\par\fi
  \ifShowChinese\par\noindent 相应修改位于#2。\par\fi
}

% 自动位置字段：重复使用 \At{页码}{行号或行号范围} 即可添加多个位置。
% 同一页出现多个 \At 时会自动合并；多个页面会自动补充 and / 以及。
\ExplSyntaxOn
\seq_new:N  \l__rr_location_pages_seq
\seq_new:N  \l__rr_location_lines_seq
\prop_new:N \l__rr_location_lines_prop
\bool_new:N \l__rr_location_singular_line_bool
\int_new:N  \l__rr_location_page_count_int
\int_new:N  \l__rr_location_line_count_int

\cs_generate_variant:Nn \seq_set_split:Nnn { NnV }

\cs_new_protected:Npn \__rr_location_collect:nn #1#2
  {
    \prop_get:NnNTF \l__rr_location_lines_prop {#1} \l_tmpa_tl
      {
        \prop_put:Nnx \l__rr_location_lines_prop {#1}
          { \exp_not:V \l_tmpa_tl ; \exp_not:n {#2} }
      }
      {
        \seq_put_right:Nn \l__rr_location_pages_seq {#1}
        \prop_put:Nnn \l__rr_location_lines_prop {#1} {#2}
      }
  }

\cs_new_protected:Npn \__rr_location_prepare:n #1
  {
    \seq_clear:N \l__rr_location_pages_seq
    \prop_clear:N \l__rr_location_lines_prop
    \cs_set_protected:Npn \At ##1##2
      { \__rr_location_collect:nn {##1} {##2} }
    #1
  }

\cs_new_protected:Npn \__rr_location_set_lines_for_page:n #1
  {
    \prop_get:NnN \l__rr_location_lines_prop {#1} \l_tmpa_tl
    \seq_set_split:NnV \l__rr_location_lines_seq {;} \l_tmpa_tl
    \seq_set_map:NNn \l__rr_location_lines_seq \l__rr_location_lines_seq
      { \tl_trim_spaces:n {##1} }
  }

\cs_new_protected:Npn \__rr_location_set_line_number:n #1
  {
    \__rr_location_set_lines_for_page:n {#1}
    \bool_set_false:N \l__rr_location_singular_line_bool
    \int_compare:nNnT { \seq_count:N \l__rr_location_lines_seq } = {1}
      {
        \tl_set:Nx \l_tmpa_tl
          { \seq_item:Nn \l__rr_location_lines_seq {1} }
        \tl_if_in:VnF \l_tmpa_tl {--}
          {
            \tl_if_in:VnF \l_tmpa_tl {,}
              { \bool_set_true:N \l__rr_location_singular_line_bool }
          }
      }
  }

\cs_new_protected:Npn \__rr_location_en_line_item:n #1
  {
    \tl_set:Nn \l_tmpa_tl {#1}
    \regex_replace_all:nnN { \s* , \s* } { ~and~ } \l_tmpa_tl
    \tl_use:N \l_tmpa_tl
  }

\cs_new_protected:Npn \__rr_location_en_lines:
  {
    \int_set:Nn \l__rr_location_line_count_int
      { \seq_count:N \l__rr_location_lines_seq }
    \seq_map_indexed_inline:Nn \l__rr_location_lines_seq
      {
        \int_compare:nNnT {##1} > {1}
          {
            \int_compare:nNnTF {\l__rr_location_line_count_int} = {2}
              { ~and~ }
              {
                \int_compare:nNnTF {##1} = {\l__rr_location_line_count_int}
                  { ,~and~ }
                  { ,~ }
              }
          }
        \__rr_location_en_line_item:n {##2}
      }
  }

\cs_new_protected:Npn \__rr_location_en_page:n #1
  {
    \__rr_location_set_line_number:n {#1}
    \tl_if_in:nnTF {#1} {--} {pages} {page}\space #1
    \int_compare:nNnTF {\l__rr_location_page_count_int} = {1}
      {
        ,\space
        \bool_if:NTF \l__rr_location_singular_line_bool {line} {lines}
        \space\__rr_location_en_lines:
      }
      {
        \space(
        \bool_if:NTF \l__rr_location_singular_line_bool {line} {lines}
        \space\__rr_location_en_lines:)
      }
  }

\cs_new_protected:Npn \__rr_location_en_pages:
  {
    \int_set:Nn \l__rr_location_page_count_int
      { \seq_count:N \l__rr_location_pages_seq }
    \seq_map_indexed_inline:Nn \l__rr_location_pages_seq
      {
        \int_compare:nNnT {##1} > {1}
          {
            \int_compare:nNnTF {\l__rr_location_page_count_int} = {2}
              { \space and\space }
              {
                \int_compare:nNnTF {##1} = {\l__rr_location_page_count_int}
                  { ,\space and\space }
                  { ,\space }
              }
          }
        \__rr_location_en_page:n {##2}
      }
  }

\cs_new_protected:Npn \__rr_location_cn_lines:
  {
    \int_set:Nn \l__rr_location_line_count_int
      { \seq_count:N \l__rr_location_lines_seq }
    \seq_map_indexed_inline:Nn \l__rr_location_lines_seq
      {
        \int_compare:nNnT {##1} > {1}
          {
            \int_compare:nNnTF {##1} = {\l__rr_location_line_count_int}
              { 和 }
              { 、 }
          }
        第~##2~行
      }
  }

\cs_new_protected:Npn \__rr_location_cn_page:n #1
  {
    \__rr_location_set_lines_for_page:n {#1}
    第~#1~页
    \int_compare:nNnTF {\l__rr_location_page_count_int} = {1}
      { \__rr_location_cn_lines: }
      { （\__rr_location_cn_lines:） }
  }

\cs_new_protected:Npn \__rr_location_cn_pages:
  {
    \int_set:Nn \l__rr_location_page_count_int
      { \seq_count:N \l__rr_location_pages_seq }
    \seq_map_indexed_inline:Nn \l__rr_location_pages_seq
      {
        \int_compare:nNnT {##1} > {1}
          {
            \int_compare:nNnTF {\l__rr_location_page_count_int} = {2}
              { 以及 }
              {
                \int_compare:nNnTF {##1} = {\l__rr_location_page_count_int}
                  { 以及 }
                  { 、 }
              }
          }
        \__rr_location_cn_page:n {##2}
      }
  }

\cs_new_protected:Npn \__rr_render_auto_location:n #1
  {
    \__rr_location_prepare:n {#1}
    \ifShowEnglish
      \par\noindent
      \int_compare:nNnTF { \seq_count:N \l__rr_location_pages_seq } = {1}
        { The~corresponding~change~appears~in~the~revised~manuscript~on~ }
        { The~corresponding~changes~appear~in~the~revised~manuscript~on~ }
      \__rr_location_en_pages:.\par
    \fi
    \ifShowChinese
      \par\noindent
      相应修改位于修订稿
      \__rr_location_cn_pages:。\par
    \fi
  }
\NewDocumentCommand{\RenderAutoLocation}{+m}{%
  \__rr_render_auto_location:n {#1}%
}
\ExplSyntaxOff

% -----------------------------------------------------------------------------
% 一条审稿意见的友好字段环境
%
% \begin{ReviewItem}{稳定编号}
%   \CommentEN{完整英文意见}
%   \CommentCN{完整中文意见}
%   \ResponseEN{英文固定首句之后的回复正文}
%   \ResponseCN{中文固定首句之后的回复正文}
%   \RevisedEN{可选：英文修订原文}
%   \RevisedCN{可选：中文对照}
%   \Evidence{可选：图、表或公式}
%   \Location{\At{8}{210--225}\At{8}{244--260}\At{10}{301}\At{12}{365}}
%   特殊位置无法用页码和行号表达时，改用 \LocationEN 和 \LocationCN 手写。
% \end{ReviewItem}
% -----------------------------------------------------------------------------
\NewDocumentEnvironment{ReviewItem}{m +b}{%
  \begingroup
  \long\def\RRCommentEN{xxx}%
  \long\def\RRCommentCN{xxx}%
  \long\def\RRResponseEN{xxx}%
  \long\def\RRResponseCN{xxx}%
  \long\def\RRRevisedEN{}%
  \long\def\RRRevisedCN{}%
  \long\def\RREvidence{}%
  \long\def\RRAutoLocation{}%
  \long\def\RRLocationEN{xxx}%
  \long\def\RRLocationCN{xxx}%

  \long\def\CommentEN##1{\long\def\RRCommentEN{##1}}%
  \long\def\CommentCN##1{\long\def\RRCommentCN{##1}}%
  \long\def\ResponseEN##1{\long\def\RRResponseEN{##1}}%
  \long\def\ResponseCN##1{\long\def\RRResponseCN{##1}}%
  \long\def\RevisedEN##1{\long\def\RRRevisedEN{##1}}%
  \long\def\RevisedCN##1{\long\def\RRRevisedCN{##1}}%
  \long\def\Evidence##1{\long\def\RREvidence{##1}}%
  \long\def\Location##1{\long\def\RRAutoLocation{##1}}%
  \long\def\LocationEN##1{\long\def\RRLocationEN{##1}}%
  \long\def\LocationCN##1{\long\def\RRLocationCN{##1}}%

  #2%

  \RenderCommentBlock{#1}{\RRCommentEN}{\RRCommentCN}%
  \RenderResponseBlock{\RRResponseEN}{\RRResponseCN}%
  \ifdefempty{\RRRevisedEN}{%
    \ifdefempty{\RRRevisedCN}{}{\RenderRevisedText{\RRRevisedEN}{\RRRevisedCN}}%
  }{\RenderRevisedText{\RRRevisedEN}{\RRRevisedCN}}%
  \ifdefempty{\RREvidence}{}{\RREvidence}%
  \ifdefempty{\RRAutoLocation}{%
    \RenderLocation{\RRLocationEN}{\RRLocationCN}%
  }{\RenderAutoLocation{\RRAutoLocation}}%
  \endgroup
}{}

\NewDocumentCommand{\FigurePlaceholder}{O{0.66\linewidth} m m}{%
  \begin{figure}[htbp]
    \centering
    \fbox{\rule{0pt}{30mm}\rule{#1}{0pt}}
    \caption{\ModeInline{#2}{#3}}
  \end{figure}
}

\newcolumntype{Y}{>{\raggedright\arraybackslash}X}

\begin{document}

\MakeLetterHeader

% =============================================================================
%                      用户编辑区 2：逐条审稿意见与回复
% 每个字段均填写一整段或多段内容，不要逐句拆分。
% =============================================================================

\ReviewerSection{1}

\begin{ReviewItem}{R1.1}
  % 审稿意见英文原文（完整粘贴，不要改写）
  \CommentEN{xxx}
  % 审稿意见中文译文（完整填写）
  \CommentCN{xxx}
  % 英文回复正文；固定首句由模板自动添加，只需填写首句之后的内容
  \ResponseEN{%
    Specifically, xxx. Accordingly, we have revised xxx to improve xxx.%
  }
  % 中文回复正文；固定首句由模板自动添加，只需填写首句之后的内容
  \ResponseCN{%
    具体而言，xxx。因此，我们已对 xxx 进行了修改，以进一步完善 xxx。%
  }
  % 可选：修订稿中的英文原文；不需要引用时可删除本行
  \RevisedEN{xxx}
  % 可选：修订原文的中文对照；不需要引用时可删除本行
  \RevisedCN{xxx}
  % 修改位置：每个 \At{页码}{行号} 表示一处；重复 \At 即可添加多个位置
  % 单个位置示例：\Location{\At{8}{210--225}}
  % 多个位置示例（相同页码会自动合并）：
  % \Location{\At{8}{210--225}\At{8}{244--260}\At{10}{301}\At{12}{365}}
  \Location{\At{xxx}{xxx}}
\end{ReviewItem}

\begin{ReviewItem}{R1.2}
  % 审稿意见英文原文（完整粘贴，不要改写）
  \CommentEN{xxx}
  % 审稿意见中文译文（完整填写）
  \CommentCN{xxx}
  % 英文回复正文；固定首句由模板自动添加
  \ResponseEN{%
    We agree that the original presentation did not make this point sufficiently
    clear. We have therefore clarified xxx and added xxx to support this revision.%
  }
  % 中文回复正文；固定首句由模板自动添加
  \ResponseCN{%
    我们同意原稿对此说明得不够清楚。因此，我们已进一步澄清 xxx，并补充 xxx 作为
    支撑。%
  }
  % 可选：回复中需要展示的图、表或公式；不需要时可删除整个 Evidence 字段
  \Evidence{%
    % 替换为真实图片时使用：
    % \includegraphics[width=0.66\linewidth]{figures/your-figure.pdf}
    \FigurePlaceholder
      {xxx}
      {xxx}%
  }
  % 修改位置：重复 \At{页码}{行号或行号范围} 可添加多个位置
  \Location{\At{xxx}{xxx}}
\end{ReviewItem}

\begin{ReviewItem}{R1.3}
  % 审稿意见英文原文（完整粘贴，不要改写）
  \CommentEN{xxx}
  % 审稿意见中文译文（完整填写）
  \CommentCN{xxx}
  % 英文回复正文；固定首句由模板自动添加
  \ResponseEN{%
    The requested comparison has been added and is summarized in
    Table~\ref{tab:comparison}.%
  }
  % 中文回复正文；固定首句由模板自动添加
  \ResponseCN{%
    我们已补充所要求的比较，结果汇总于表~\ref{tab:comparison}。%
  }
  % 可选：定量比较表；将表题、方法名、指标名和数值中的 xxx 替换为实际内容
  \Evidence{%
    \begin{table}[htbp]
      \centering
      \caption{\ModeInline
        {xxx}
        {xxx}}
      \label{tab:comparison}
      \begin{tabularx}{0.86\linewidth}{@{}Y>{\centering\arraybackslash}p{2.6cm}>{\centering\arraybackslash}p{2.6cm}@{}}
        \toprule
        \ModeInline{Method}{方法} & \ModeInline{Metric 1}{指标 1} & \ModeInline{Metric 2}{指标 2} \\
        \midrule
        \ModeInline{Proposed method}{本文方法} & xxx & xxx \\
        \ModeInline{Baseline method}{基线方法} & xxx & xxx \\
        \bottomrule
      \end{tabularx}
    \end{table}%
  }
  % 修改位置：重复 \At{页码}{行号或行号范围} 可添加多个位置
  \Location{\At{xxx}{xxx}}
\end{ReviewItem}

\ReviewerSection{2}

\begin{ReviewItem}{R2.1}
  % 审稿意见英文原文（完整粘贴，不要改写）
  \CommentEN{xxx}
  % 审稿意见中文译文（完整填写）
  \CommentCN{xxx}
  % 英文回复正文；固定首句由模板自动添加
  \ResponseEN{%
    In response, we have revised xxx and added xxx to substantiate this point.%
  }
  % 中文回复正文；固定首句由模板自动添加
  \ResponseCN{%
    作为回应，我们已修改 xxx，并补充 xxx 以提供支撑。%
  }
  % 修改位置：重复 \At{页码}{行号或行号范围} 可添加多个位置
  \Location{\At{xxx}{xxx}}
\end{ReviewItem}

% 添加新意见：复制一个完整的 ReviewItem 环境。
% 添加新审稿人：复制 ReviewerSection，并从新的 ReviewItem 开始。

\end{document}
"
% =============================================================================
\providecommand{\ResponseMode}{both}

\newif\ifShowEnglish
\newif\ifShowChinese

\newcommand{\ResponseDisplayMode}[1]{%
  \ShowEnglishfalse
  \ShowChinesefalse
  \ifstrequal{#1}{en}{%
    \ShowEnglishtrue
  }{%
    \ifstrequal{#1}{cn}{%
      \ShowChinesetrue
    }{%
      \ifstrequal{#1}{both}{%
        \ShowEnglishtrue
        \ShowChinesetrue
      }{%
        \PackageError{ReviewResponse}{Unknown display mode '#1'}%
          {Use en, cn, or both.}%
      }%
    }%
  }%
}
\expandafter\ResponseDisplayMode\expandafter{\ResponseMode}

\newcommand{\ModeInline}[2]{%
  \ifShowEnglish #1\ifShowChinese\ /\ \fi\fi
  \ifShowChinese #2\fi
}

\renewcommand{\figurename}{\ModeInline{Figure}{图}}
\renewcommand{\tablename}{\ModeInline{Table}{表}}

% =============================================================================
%                         用户编辑区 1：信件整体字段
% 只需在这里填写一次，标题、开场说明和署名会自动同步替换。
% 未确定的内容暂时保留为普通黑色 xxx。
% =============================================================================

% 论文英文标题
\newcommand{\MetaTitleEN}{xxx}
% 论文中文标题
\newcommand{\MetaTitleCN}{xxx}
% 稿件编号（中英文共用，只填一次）
\newcommand{\MetaManuscriptID}{xxx}
% 通讯作者英文姓名
\newcommand{\MetaCorrespondingAuthorEN}{xxx}
% 通讯作者中文姓名
\newcommand{\MetaCorrespondingAuthorCN}{xxx}

\newcommand{\FieldSalutationEN}{Dear Editor and Reviewers,}
\newcommand{\FieldSalutationCN}{尊敬的编辑和审稿专家：}

% 每条回复的固定首句，由模板自动添加，无需在 ResponseEN / ResponseCN 中重复填写。
\newcommand{\ResponseOpeningEN}{We thank the reviewer for these comments.}
\newcommand{\ResponseOpeningCN}{感谢审稿人提出这些意见。}

\newcommand{\FieldOpeningEN}{%
  Thank you for your careful evaluation and constructive comments. We have
  carefully considered each comment and revised the manuscript accordingly. Our
  point-by-point responses are provided below. The reviewers' comments are
  reproduced in bold, followed by our responses. All changes in the revised
  manuscript are highlighted in red.%
}

\newcommand{\FieldOpeningCN}{%
  感谢您对本稿件的认真评阅和建设性意见。我们已逐条认真研究并据此修订稿件。下文
  给出逐点回复。审稿意见以粗体显示，随后给出我们的回复。修订稿中的所有改动均以
  红色标出。%
}

\newcommand{\FieldSignOffEN}{%
  With kindest regards,\par
  Yours Sincerely,\par
  \MetaCorrespondingAuthorEN%
}

\newcommand{\FieldSignOffCN}{%
  此致\par
  敬礼！\par
  \MetaCorrespondingAuthorCN%
}

% =============================================================================
%                         以下为排版逻辑，通常无需修改
% =============================================================================

\newcommand{\IndentedParagraph}[1]{%
  \par\begingroup
  \setlength{\parindent}{21pt}%
  \indent #1\par
  \endgroup
}

\newcommand{\MakeLetterHeader}{%
  \begin{center}
    \bfseries\fontsize{12pt}{14.4pt}\selectfont
    \ifShowEnglish Revisions to the paper titled \MetaTitleEN\fi
    \ifShowEnglish\ifShowChinese\\[0.35em]\fi\fi
    \ifShowChinese 关于论文《\MetaTitleCN》的修改说明\fi
  \end{center}
  \vspace{1.2\baselineskip}

  \ifShowEnglish\noindent Manuscript Number: \MetaManuscriptID\par\fi
  \ifShowChinese\noindent 稿件编号：\MetaManuscriptID\par\fi
  \vspace{1.9\baselineskip}

  \ifShowEnglish \FieldSalutationEN\par\fi
  \ifShowEnglish\ifShowChinese\smallskip\fi\fi
  \ifShowChinese \FieldSalutationCN\par\fi
  \medskip

  \ifShowEnglish\IndentedParagraph{\FieldOpeningEN}\fi
  \ifShowEnglish\ifShowChinese\vspace{0.45\baselineskip}\fi\fi
  \ifShowChinese\IndentedParagraph{\FieldOpeningCN}\fi

  \vspace{0.65\baselineskip}
  \ifShowEnglish \FieldSignOffEN\par\fi
  \ifShowEnglish\ifShowChinese\vspace{0.45\baselineskip}\fi\fi
  \ifShowChinese \FieldSignOffCN\par\fi
}

\newcounter{reviewcomment}

\newcommand{\EditorSection}{%
  \vspace{1.35\baselineskip}
  \noindent\textbf{\ModeInline{Response to the Editor:}{对编辑意见的回复：}}\par
  \vspace{0.75\baselineskip}
  \setcounter{reviewcomment}{0}%
}

\newcommand{\ReviewerSection}[1]{%
  \vspace{1.35\baselineskip}
  \noindent\textbf{\ModeInline{Response to Reviewer \##1:}{对审稿人 \##1 的回复：}}\par
  \vspace{0.75\baselineskip}
  \setcounter{reviewcomment}{0}%
}

% 审稿意见：先显示完整英文块，再显示完整中文块；不做逐句交替。
\NewDocumentCommand{\RenderCommentBlock}{m +m +m}{%
  \refstepcounter{reviewcomment}%
  \Needspace{7\baselineskip}%
  \par\vspace{0.9\baselineskip}
  \noindent\hypertarget{#1}{}%
  \begingroup\bfseries
  \arabic{reviewcomment}.\ %
  \ifShowEnglish #2\fi
  \ifShowEnglish\ifShowChinese\par\smallskip\noindent\hspace*{1.35em}\fi\fi
  \ifShowChinese #3\fi
  \par\endgroup
  \nobreak\smallskip
}

% 作者回复：先显示完整英文块，再显示完整中文块；不做逐句交替。
\NewDocumentCommand{\RenderResponseBlock}{+m +m}{%
  \Needspace{4\baselineskip}%
  \noindent\textbf{\ModeInline{Response:}{回复：}}\par
  \nobreak
  \ifShowEnglish
    \noindent\mbox{\textgreater\kern0.08em\textgreater}\
    \ResponseOpeningEN\space #1\par
  \fi
  \ifShowEnglish\ifShowChinese\smallskip\fi\fi
  \ifShowChinese
    \noindent\mbox{\textgreater\kern0.08em\textgreater}\
    \ResponseOpeningCN\space #2\par
  \fi
}

\NewDocumentCommand{\RenderRevisedText}{+m +m}{%
  \ifShowEnglish
    \par\begingroup
    \setlength{\leftskip}{21pt}\setlength{\rightskip}{21pt}%
    \textit{Revised text: }``#1''\par
    \endgroup
  \fi
  \ifShowEnglish\ifShowChinese\smallskip\fi\fi
  \ifShowChinese
    \par\begingroup
    \setlength{\leftskip}{21pt}\setlength{\rightskip}{21pt}%
    修订文本：“#2”\par
    \endgroup
  \fi
}

\NewDocumentCommand{\RenderLocation}{+m +m}{%
  \ifShowEnglish\par\noindent The corresponding change appears in #1.\par\fi
  \ifShowChinese\par\noindent 相应修改位于#2。\par\fi
}

% 自动位置字段：重复使用 \At{页码}{行号或行号范围} 即可添加多个位置。
% 同一页出现多个 \At 时会自动合并；多个页面会自动补充 and / 以及。
\ExplSyntaxOn
\seq_new:N  \l__rr_location_pages_seq
\seq_new:N  \l__rr_location_lines_seq
\prop_new:N \l__rr_location_lines_prop
\bool_new:N \l__rr_location_singular_line_bool
\int_new:N  \l__rr_location_page_count_int
\int_new:N  \l__rr_location_line_count_int

\cs_generate_variant:Nn \seq_set_split:Nnn { NnV }

\cs_new_protected:Npn \__rr_location_collect:nn #1#2
  {
    \prop_get:NnNTF \l__rr_location_lines_prop {#1} \l_tmpa_tl
      {
        \prop_put:Nnx \l__rr_location_lines_prop {#1}
          { \exp_not:V \l_tmpa_tl ; \exp_not:n {#2} }
      }
      {
        \seq_put_right:Nn \l__rr_location_pages_seq {#1}
        \prop_put:Nnn \l__rr_location_lines_prop {#1} {#2}
      }
  }

\cs_new_protected:Npn \__rr_location_prepare:n #1
  {
    \seq_clear:N \l__rr_location_pages_seq
    \prop_clear:N \l__rr_location_lines_prop
    \cs_set_protected:Npn \At ##1##2
      { \__rr_location_collect:nn {##1} {##2} }
    #1
  }

\cs_new_protected:Npn \__rr_location_set_lines_for_page:n #1
  {
    \prop_get:NnN \l__rr_location_lines_prop {#1} \l_tmpa_tl
    \seq_set_split:NnV \l__rr_location_lines_seq {;} \l_tmpa_tl
    \seq_set_map:NNn \l__rr_location_lines_seq \l__rr_location_lines_seq
      { \tl_trim_spaces:n {##1} }
  }

\cs_new_protected:Npn \__rr_location_set_line_number:n #1
  {
    \__rr_location_set_lines_for_page:n {#1}
    \bool_set_false:N \l__rr_location_singular_line_bool
    \int_compare:nNnT { \seq_count:N \l__rr_location_lines_seq } = {1}
      {
        \tl_set:Nx \l_tmpa_tl
          { \seq_item:Nn \l__rr_location_lines_seq {1} }
        \tl_if_in:VnF \l_tmpa_tl {--}
          {
            \tl_if_in:VnF \l_tmpa_tl {,}
              { \bool_set_true:N \l__rr_location_singular_line_bool }
          }
      }
  }

\cs_new_protected:Npn \__rr_location_en_line_item:n #1
  {
    \tl_set:Nn \l_tmpa_tl {#1}
    \regex_replace_all:nnN { \s* , \s* } { ~and~ } \l_tmpa_tl
    \tl_use:N \l_tmpa_tl
  }

\cs_new_protected:Npn \__rr_location_en_lines:
  {
    \int_set:Nn \l__rr_location_line_count_int
      { \seq_count:N \l__rr_location_lines_seq }
    \seq_map_indexed_inline:Nn \l__rr_location_lines_seq
      {
        \int_compare:nNnT {##1} > {1}
          {
            \int_compare:nNnTF {\l__rr_location_line_count_int} = {2}
              { ~and~ }
              {
                \int_compare:nNnTF {##1} = {\l__rr_location_line_count_int}
                  { ,~and~ }
                  { ,~ }
              }
          }
        \__rr_location_en_line_item:n {##2}
      }
  }

\cs_new_protected:Npn \__rr_location_en_page:n #1
  {
    \__rr_location_set_line_number:n {#1}
    \tl_if_in:nnTF {#1} {--} {pages} {page}\space #1
    \int_compare:nNnTF {\l__rr_location_page_count_int} = {1}
      {
        ,\space
        \bool_if:NTF \l__rr_location_singular_line_bool {line} {lines}
        \space\__rr_location_en_lines:
      }
      {
        \space(
        \bool_if:NTF \l__rr_location_singular_line_bool {line} {lines}
        \space\__rr_location_en_lines:)
      }
  }

\cs_new_protected:Npn \__rr_location_en_pages:
  {
    \int_set:Nn \l__rr_location_page_count_int
      { \seq_count:N \l__rr_location_pages_seq }
    \seq_map_indexed_inline:Nn \l__rr_location_pages_seq
      {
        \int_compare:nNnT {##1} > {1}
          {
            \int_compare:nNnTF {\l__rr_location_page_count_int} = {2}
              { \space and\space }
              {
                \int_compare:nNnTF {##1} = {\l__rr_location_page_count_int}
                  { ,\space and\space }
                  { ,\space }
              }
          }
        \__rr_location_en_page:n {##2}
      }
  }

\cs_new_protected:Npn \__rr_location_cn_lines:
  {
    \int_set:Nn \l__rr_location_line_count_int
      { \seq_count:N \l__rr_location_lines_seq }
    \seq_map_indexed_inline:Nn \l__rr_location_lines_seq
      {
        \int_compare:nNnT {##1} > {1}
          {
            \int_compare:nNnTF {##1} = {\l__rr_location_line_count_int}
              { 和 }
              { 、 }
          }
        第~##2~行
      }
  }

\cs_new_protected:Npn \__rr_location_cn_page:n #1
  {
    \__rr_location_set_lines_for_page:n {#1}
    第~#1~页
    \int_compare:nNnTF {\l__rr_location_page_count_int} = {1}
      { \__rr_location_cn_lines: }
      { （\__rr_location_cn_lines:） }
  }

\cs_new_protected:Npn \__rr_location_cn_pages:
  {
    \int_set:Nn \l__rr_location_page_count_int
      { \seq_count:N \l__rr_location_pages_seq }
    \seq_map_indexed_inline:Nn \l__rr_location_pages_seq
      {
        \int_compare:nNnT {##1} > {1}
          {
            \int_compare:nNnTF {\l__rr_location_page_count_int} = {2}
              { 以及 }
              {
                \int_compare:nNnTF {##1} = {\l__rr_location_page_count_int}
                  { 以及 }
                  { 、 }
              }
          }
        \__rr_location_cn_page:n {##2}
      }
  }

\cs_new_protected:Npn \__rr_render_auto_location:n #1
  {
    \__rr_location_prepare:n {#1}
    \ifShowEnglish
      \par\noindent
      \int_compare:nNnTF { \seq_count:N \l__rr_location_pages_seq } = {1}
        { The~corresponding~change~appears~in~the~revised~manuscript~on~ }
        { The~corresponding~changes~appear~in~the~revised~manuscript~on~ }
      \__rr_location_en_pages:.\par
    \fi
    \ifShowChinese
      \par\noindent
      相应修改位于修订稿
      \__rr_location_cn_pages:。\par
    \fi
  }
\NewDocumentCommand{\RenderAutoLocation}{+m}{%
  \__rr_render_auto_location:n {#1}%
}
\ExplSyntaxOff

% -----------------------------------------------------------------------------
% 一条审稿意见的友好字段环境
%
% \begin{ReviewItem}{稳定编号}
%   \CommentEN{完整英文意见}
%   \CommentCN{完整中文意见}
%   \ResponseEN{英文固定首句之后的回复正文}
%   \ResponseCN{中文固定首句之后的回复正文}
%   \RevisedEN{可选：英文修订原文}
%   \RevisedCN{可选：中文对照}
%   \Evidence{可选：图、表或公式}
%   \Location{\At{8}{210--225}\At{8}{244--260}\At{10}{301}\At{12}{365}}
%   特殊位置无法用页码和行号表达时，改用 \LocationEN 和 \LocationCN 手写。
% \end{ReviewItem}
% -----------------------------------------------------------------------------
\NewDocumentEnvironment{ReviewItem}{m +b}{%
  \begingroup
  \long\def\RRCommentEN{xxx}%
  \long\def\RRCommentCN{xxx}%
  \long\def\RRResponseEN{xxx}%
  \long\def\RRResponseCN{xxx}%
  \long\def\RRRevisedEN{}%
  \long\def\RRRevisedCN{}%
  \long\def\RREvidence{}%
  \long\def\RRAutoLocation{}%
  \long\def\RRLocationEN{xxx}%
  \long\def\RRLocationCN{xxx}%

  \long\def\CommentEN##1{\long\def\RRCommentEN{##1}}%
  \long\def\CommentCN##1{\long\def\RRCommentCN{##1}}%
  \long\def\ResponseEN##1{\long\def\RRResponseEN{##1}}%
  \long\def\ResponseCN##1{\long\def\RRResponseCN{##1}}%
  \long\def\RevisedEN##1{\long\def\RRRevisedEN{##1}}%
  \long\def\RevisedCN##1{\long\def\RRRevisedCN{##1}}%
  \long\def\Evidence##1{\long\def\RREvidence{##1}}%
  \long\def\Location##1{\long\def\RRAutoLocation{##1}}%
  \long\def\LocationEN##1{\long\def\RRLocationEN{##1}}%
  \long\def\LocationCN##1{\long\def\RRLocationCN{##1}}%

  #2%

  \RenderCommentBlock{#1}{\RRCommentEN}{\RRCommentCN}%
  \RenderResponseBlock{\RRResponseEN}{\RRResponseCN}%
  \ifdefempty{\RRRevisedEN}{%
    \ifdefempty{\RRRevisedCN}{}{\RenderRevisedText{\RRRevisedEN}{\RRRevisedCN}}%
  }{\RenderRevisedText{\RRRevisedEN}{\RRRevisedCN}}%
  \ifdefempty{\RREvidence}{}{\RREvidence}%
  \ifdefempty{\RRAutoLocation}{%
    \RenderLocation{\RRLocationEN}{\RRLocationCN}%
  }{\RenderAutoLocation{\RRAutoLocation}}%
  \endgroup
}{}

\NewDocumentCommand{\FigurePlaceholder}{O{0.66\linewidth} m m}{%
  \begin{figure}[htbp]
    \centering
    \fbox{\rule{0pt}{30mm}\rule{#1}{0pt}}
    \caption{\ModeInline{#2}{#3}}
  \end{figure}
}

\newcolumntype{Y}{>{\raggedright\arraybackslash}X}

\begin{document}

\MakeLetterHeader

% =============================================================================
%                      用户编辑区 2：逐条审稿意见与回复
% 每个字段均填写一整段或多段内容，不要逐句拆分。
% =============================================================================

\ReviewerSection{1}

\begin{ReviewItem}{R1.1}
  % 审稿意见英文原文（完整粘贴，不要改写）
  \CommentEN{xxx}
  % 审稿意见中文译文（完整填写）
  \CommentCN{xxx}
  % 英文回复正文；固定首句由模板自动添加，只需填写首句之后的内容
  \ResponseEN{%
    Specifically, xxx. Accordingly, we have revised xxx to improve xxx.%
  }
  % 中文回复正文；固定首句由模板自动添加，只需填写首句之后的内容
  \ResponseCN{%
    具体而言，xxx。因此，我们已对 xxx 进行了修改，以进一步完善 xxx。%
  }
  % 可选：修订稿中的英文原文；不需要引用时可删除本行
  \RevisedEN{xxx}
  % 可选：修订原文的中文对照；不需要引用时可删除本行
  \RevisedCN{xxx}
  % 修改位置：每个 \At{页码}{行号} 表示一处；重复 \At 即可添加多个位置
  % 单个位置示例：\Location{\At{8}{210--225}}
  % 多个位置示例（相同页码会自动合并）：
  % \Location{\At{8}{210--225}\At{8}{244--260}\At{10}{301}\At{12}{365}}
  \Location{\At{xxx}{xxx}}
\end{ReviewItem}

\begin{ReviewItem}{R1.2}
  % 审稿意见英文原文（完整粘贴，不要改写）
  \CommentEN{xxx}
  % 审稿意见中文译文（完整填写）
  \CommentCN{xxx}
  % 英文回复正文；固定首句由模板自动添加
  \ResponseEN{%
    We agree that the original presentation did not make this point sufficiently
    clear. We have therefore clarified xxx and added xxx to support this revision.%
  }
  % 中文回复正文；固定首句由模板自动添加
  \ResponseCN{%
    我们同意原稿对此说明得不够清楚。因此，我们已进一步澄清 xxx，并补充 xxx 作为
    支撑。%
  }
  % 可选：回复中需要展示的图、表或公式；不需要时可删除整个 Evidence 字段
  \Evidence{%
    % 替换为真实图片时使用：
    % \includegraphics[width=0.66\linewidth]{figures/your-figure.pdf}
    \FigurePlaceholder
      {xxx}
      {xxx}%
  }
  % 修改位置：重复 \At{页码}{行号或行号范围} 可添加多个位置
  \Location{\At{xxx}{xxx}}
\end{ReviewItem}

\begin{ReviewItem}{R1.3}
  % 审稿意见英文原文（完整粘贴，不要改写）
  \CommentEN{xxx}
  % 审稿意见中文译文（完整填写）
  \CommentCN{xxx}
  % 英文回复正文；固定首句由模板自动添加
  \ResponseEN{%
    The requested comparison has been added and is summarized in
    Table~\ref{tab:comparison}.%
  }
  % 中文回复正文；固定首句由模板自动添加
  \ResponseCN{%
    我们已补充所要求的比较，结果汇总于表~\ref{tab:comparison}。%
  }
  % 可选：定量比较表；将表题、方法名、指标名和数值中的 xxx 替换为实际内容
  \Evidence{%
    \begin{table}[htbp]
      \centering
      \caption{\ModeInline
        {xxx}
        {xxx}}
      \label{tab:comparison}
      \begin{tabularx}{0.86\linewidth}{@{}Y>{\centering\arraybackslash}p{2.6cm}>{\centering\arraybackslash}p{2.6cm}@{}}
        \toprule
        \ModeInline{Method}{方法} & \ModeInline{Metric 1}{指标 1} & \ModeInline{Metric 2}{指标 2} \\
        \midrule
        \ModeInline{Proposed method}{本文方法} & xxx & xxx \\
        \ModeInline{Baseline method}{基线方法} & xxx & xxx \\
        \bottomrule
      \end{tabularx}
    \end{table}%
  }
  % 修改位置：重复 \At{页码}{行号或行号范围} 可添加多个位置
  \Location{\At{xxx}{xxx}}
\end{ReviewItem}

\ReviewerSection{2}

\begin{ReviewItem}{R2.1}
  % 审稿意见英文原文（完整粘贴，不要改写）
  \CommentEN{xxx}
  % 审稿意见中文译文（完整填写）
  \CommentCN{xxx}
  % 英文回复正文；固定首句由模板自动添加
  \ResponseEN{%
    In response, we have revised xxx and added xxx to substantiate this point.%
  }
  % 中文回复正文；固定首句由模板自动添加
  \ResponseCN{%
    作为回应，我们已修改 xxx，并补充 xxx 以提供支撑。%
  }
  % 修改位置：重复 \At{页码}{行号或行号范围} 可添加多个位置
  \Location{\At{xxx}{xxx}}
\end{ReviewItem}

% 添加新意见：复制一个完整的 ReviewItem 环境。
% 添加新审稿人：复制 ReviewerSection，并从新的 ReviewItem 开始。

\end{document}
"
% =============================================================================
\providecommand{\ResponseMode}{both}

\newif\ifShowEnglish
\newif\ifShowChinese

\newcommand{\ResponseDisplayMode}[1]{%
  \ShowEnglishfalse
  \ShowChinesefalse
  \ifstrequal{#1}{en}{%
    \ShowEnglishtrue
  }{%
    \ifstrequal{#1}{cn}{%
      \ShowChinesetrue
    }{%
      \ifstrequal{#1}{both}{%
        \ShowEnglishtrue
        \ShowChinesetrue
      }{%
        \PackageError{ReviewResponse}{Unknown display mode '#1'}%
          {Use en, cn, or both.}%
      }%
    }%
  }%
}
\expandafter\ResponseDisplayMode\expandafter{\ResponseMode}

\newcommand{\ModeInline}[2]{%
  \ifShowEnglish #1\ifShowChinese\ /\ \fi\fi
  \ifShowChinese #2\fi
}

\renewcommand{\figurename}{\ModeInline{Figure}{图}}
\renewcommand{\tablename}{\ModeInline{Table}{表}}

% =============================================================================
%                         用户编辑区 1：信件整体字段
% 只需在这里填写一次，标题、开场说明和署名会自动同步替换。
% 未确定的内容暂时保留为普通黑色 xxx。
% =============================================================================

% 论文英文标题
\newcommand{\MetaTitleEN}{xxx}
% 论文中文标题
\newcommand{\MetaTitleCN}{xxx}
% 稿件编号（中英文共用，只填一次）
\newcommand{\MetaManuscriptID}{xxx}
% 通讯作者英文姓名
\newcommand{\MetaCorrespondingAuthorEN}{xxx}
% 通讯作者中文姓名
\newcommand{\MetaCorrespondingAuthorCN}{xxx}

\newcommand{\FieldSalutationEN}{Dear Editor and Reviewers,}
\newcommand{\FieldSalutationCN}{尊敬的编辑和审稿专家：}

% 每条回复的固定首句，由模板自动添加，无需在 ResponseEN / ResponseCN 中重复填写。
\newcommand{\ResponseOpeningEN}{We thank the reviewer for these comments.}
\newcommand{\ResponseOpeningCN}{感谢审稿人提出这些意见。}

\newcommand{\FieldOpeningEN}{%
  Thank you for your careful evaluation and constructive comments. We have
  carefully considered each comment and revised the manuscript accordingly. Our
  point-by-point responses are provided below. The reviewers' comments are
  reproduced in bold, followed by our responses. All changes in the revised
  manuscript are highlighted in red.%
}

\newcommand{\FieldOpeningCN}{%
  感谢您对本稿件的认真评阅和建设性意见。我们已逐条认真研究并据此修订稿件。下文
  给出逐点回复。审稿意见以粗体显示，随后给出我们的回复。修订稿中的所有改动均以
  红色标出。%
}

\newcommand{\FieldSignOffEN}{%
  With kindest regards,\par
  Yours Sincerely,\par
  \MetaCorrespondingAuthorEN%
}

\newcommand{\FieldSignOffCN}{%
  此致\par
  敬礼！\par
  \MetaCorrespondingAuthorCN%
}

% =============================================================================
%                         以下为排版逻辑，通常无需修改
% =============================================================================

\newcommand{\IndentedParagraph}[1]{%
  \par\begingroup
  \setlength{\parindent}{21pt}%
  \indent #1\par
  \endgroup
}

\newcommand{\MakeLetterHeader}{%
  \begin{center}
    \bfseries\fontsize{12pt}{14.4pt}\selectfont
    \ifShowEnglish Revisions to the paper titled \MetaTitleEN\fi
    \ifShowEnglish\ifShowChinese\\[0.35em]\fi\fi
    \ifShowChinese 关于论文《\MetaTitleCN》的修改说明\fi
  \end{center}
  \vspace{1.2\baselineskip}

  \ifShowEnglish\noindent Manuscript Number: \MetaManuscriptID\par\fi
  \ifShowChinese\noindent 稿件编号：\MetaManuscriptID\par\fi
  \vspace{1.9\baselineskip}

  \ifShowEnglish \FieldSalutationEN\par\fi
  \ifShowEnglish\ifShowChinese\smallskip\fi\fi
  \ifShowChinese \FieldSalutationCN\par\fi
  \medskip

  \ifShowEnglish\IndentedParagraph{\FieldOpeningEN}\fi
  \ifShowEnglish\ifShowChinese\vspace{0.45\baselineskip}\fi\fi
  \ifShowChinese\IndentedParagraph{\FieldOpeningCN}\fi

  \vspace{0.65\baselineskip}
  \ifShowEnglish \FieldSignOffEN\par\fi
  \ifShowEnglish\ifShowChinese\vspace{0.45\baselineskip}\fi\fi
  \ifShowChinese \FieldSignOffCN\par\fi
}

\newcounter{reviewcomment}

\newcommand{\EditorSection}{%
  \vspace{1.35\baselineskip}
  \noindent\textbf{\ModeInline{Response to the Editor:}{对编辑意见的回复：}}\par
  \vspace{0.75\baselineskip}
  \setcounter{reviewcomment}{0}%
}

\newcommand{\ReviewerSection}[1]{%
  \vspace{1.35\baselineskip}
  \noindent\textbf{\ModeInline{Response to Reviewer \##1:}{对审稿人 \##1 的回复：}}\par
  \vspace{0.75\baselineskip}
  \setcounter{reviewcomment}{0}%
}

% 审稿意见：先显示完整英文块，再显示完整中文块；不做逐句交替。
\NewDocumentCommand{\RenderCommentBlock}{m +m +m}{%
  \refstepcounter{reviewcomment}%
  \Needspace{7\baselineskip}%
  \par\vspace{0.9\baselineskip}
  \noindent\hypertarget{#1}{}%
  \begingroup\bfseries
  \arabic{reviewcomment}.\ %
  \ifShowEnglish #2\fi
  \ifShowEnglish\ifShowChinese\par\smallskip\noindent\hspace*{1.35em}\fi\fi
  \ifShowChinese #3\fi
  \par\endgroup
  \nobreak\smallskip
}

% 作者回复：先显示完整英文块，再显示完整中文块；不做逐句交替。
\NewDocumentCommand{\RenderResponseBlock}{+m +m}{%
  \Needspace{4\baselineskip}%
  \noindent\textbf{\ModeInline{Response:}{回复：}}\par
  \nobreak
  \ifShowEnglish
    \noindent\mbox{\textgreater\kern0.08em\textgreater}\
    \ResponseOpeningEN\space #1\par
  \fi
  \ifShowEnglish\ifShowChinese\smallskip\fi\fi
  \ifShowChinese
    \noindent\mbox{\textgreater\kern0.08em\textgreater}\
    \ResponseOpeningCN\space #2\par
  \fi
}

\NewDocumentCommand{\RenderRevisedText}{+m +m}{%
  \ifShowEnglish
    \par\begingroup
    \setlength{\leftskip}{21pt}\setlength{\rightskip}{21pt}%
    \textit{Revised text: }``#1''\par
    \endgroup
  \fi
  \ifShowEnglish\ifShowChinese\smallskip\fi\fi
  \ifShowChinese
    \par\begingroup
    \setlength{\leftskip}{21pt}\setlength{\rightskip}{21pt}%
    修订文本：“#2”\par
    \endgroup
  \fi
}

\NewDocumentCommand{\RenderLocation}{+m +m}{%
  \ifShowEnglish\par\noindent The corresponding change appears in #1.\par\fi
  \ifShowChinese\par\noindent 相应修改位于#2。\par\fi
}

% 自动位置字段：重复使用 \At{页码}{行号或行号范围} 即可添加多个位置。
% 同一页出现多个 \At 时会自动合并；多个页面会自动补充 and / 以及。
\ExplSyntaxOn
\seq_new:N  \l__rr_location_pages_seq
\seq_new:N  \l__rr_location_lines_seq
\prop_new:N \l__rr_location_lines_prop
\bool_new:N \l__rr_location_singular_line_bool
\int_new:N  \l__rr_location_page_count_int
\int_new:N  \l__rr_location_line_count_int

\cs_generate_variant:Nn \seq_set_split:Nnn { NnV }

\cs_new_protected:Npn \__rr_location_collect:nn #1#2
  {
    \prop_get:NnNTF \l__rr_location_lines_prop {#1} \l_tmpa_tl
      {
        \prop_put:Nnx \l__rr_location_lines_prop {#1}
          { \exp_not:V \l_tmpa_tl ; \exp_not:n {#2} }
      }
      {
        \seq_put_right:Nn \l__rr_location_pages_seq {#1}
        \prop_put:Nnn \l__rr_location_lines_prop {#1} {#2}
      }
  }

\cs_new_protected:Npn \__rr_location_prepare:n #1
  {
    \seq_clear:N \l__rr_location_pages_seq
    \prop_clear:N \l__rr_location_lines_prop
    \cs_set_protected:Npn \At ##1##2
      { \__rr_location_collect:nn {##1} {##2} }
    #1
  }

\cs_new_protected:Npn \__rr_location_set_lines_for_page:n #1
  {
    \prop_get:NnN \l__rr_location_lines_prop {#1} \l_tmpa_tl
    \seq_set_split:NnV \l__rr_location_lines_seq {;} \l_tmpa_tl
    \seq_set_map:NNn \l__rr_location_lines_seq \l__rr_location_lines_seq
      { \tl_trim_spaces:n {##1} }
  }

\cs_new_protected:Npn \__rr_location_set_line_number:n #1
  {
    \__rr_location_set_lines_for_page:n {#1}
    \bool_set_false:N \l__rr_location_singular_line_bool
    \int_compare:nNnT { \seq_count:N \l__rr_location_lines_seq } = {1}
      {
        \tl_set:Nx \l_tmpa_tl
          { \seq_item:Nn \l__rr_location_lines_seq {1} }
        \tl_if_in:VnF \l_tmpa_tl {--}
          {
            \tl_if_in:VnF \l_tmpa_tl {,}
              { \bool_set_true:N \l__rr_location_singular_line_bool }
          }
      }
  }

\cs_new_protected:Npn \__rr_location_en_line_item:n #1
  {
    \tl_set:Nn \l_tmpa_tl {#1}
    \regex_replace_all:nnN { \s* , \s* } { ~and~ } \l_tmpa_tl
    \tl_use:N \l_tmpa_tl
  }

\cs_new_protected:Npn \__rr_location_en_lines:
  {
    \int_set:Nn \l__rr_location_line_count_int
      { \seq_count:N \l__rr_location_lines_seq }
    \seq_map_indexed_inline:Nn \l__rr_location_lines_seq
      {
        \int_compare:nNnT {##1} > {1}
          {
            \int_compare:nNnTF {\l__rr_location_line_count_int} = {2}
              { ~and~ }
              {
                \int_compare:nNnTF {##1} = {\l__rr_location_line_count_int}
                  { ,~and~ }
                  { ,~ }
              }
          }
        \__rr_location_en_line_item:n {##2}
      }
  }

\cs_new_protected:Npn \__rr_location_en_page:n #1
  {
    \__rr_location_set_line_number:n {#1}
    \tl_if_in:nnTF {#1} {--} {pages} {page}\space #1
    \int_compare:nNnTF {\l__rr_location_page_count_int} = {1}
      {
        ,\space
        \bool_if:NTF \l__rr_location_singular_line_bool {line} {lines}
        \space\__rr_location_en_lines:
      }
      {
        \space(
        \bool_if:NTF \l__rr_location_singular_line_bool {line} {lines}
        \space\__rr_location_en_lines:)
      }
  }

\cs_new_protected:Npn \__rr_location_en_pages:
  {
    \int_set:Nn \l__rr_location_page_count_int
      { \seq_count:N \l__rr_location_pages_seq }
    \seq_map_indexed_inline:Nn \l__rr_location_pages_seq
      {
        \int_compare:nNnT {##1} > {1}
          {
            \int_compare:nNnTF {\l__rr_location_page_count_int} = {2}
              { \space and\space }
              {
                \int_compare:nNnTF {##1} = {\l__rr_location_page_count_int}
                  { ,\space and\space }
                  { ,\space }
              }
          }
        \__rr_location_en_page:n {##2}
      }
  }

\cs_new_protected:Npn \__rr_location_cn_lines:
  {
    \int_set:Nn \l__rr_location_line_count_int
      { \seq_count:N \l__rr_location_lines_seq }
    \seq_map_indexed_inline:Nn \l__rr_location_lines_seq
      {
        \int_compare:nNnT {##1} > {1}
          {
            \int_compare:nNnTF {##1} = {\l__rr_location_line_count_int}
              { 和 }
              { 、 }
          }
        第~##2~行
      }
  }

\cs_new_protected:Npn \__rr_location_cn_page:n #1
  {
    \__rr_location_set_lines_for_page:n {#1}
    第~#1~页
    \int_compare:nNnTF {\l__rr_location_page_count_int} = {1}
      { \__rr_location_cn_lines: }
      { （\__rr_location_cn_lines:） }
  }

\cs_new_protected:Npn \__rr_location_cn_pages:
  {
    \int_set:Nn \l__rr_location_page_count_int
      { \seq_count:N \l__rr_location_pages_seq }
    \seq_map_indexed_inline:Nn \l__rr_location_pages_seq
      {
        \int_compare:nNnT {##1} > {1}
          {
            \int_compare:nNnTF {\l__rr_location_page_count_int} = {2}
              { 以及 }
              {
                \int_compare:nNnTF {##1} = {\l__rr_location_page_count_int}
                  { 以及 }
                  { 、 }
              }
          }
        \__rr_location_cn_page:n {##2}
      }
  }

\cs_new_protected:Npn \__rr_render_auto_location:n #1
  {
    \__rr_location_prepare:n {#1}
    \ifShowEnglish
      \par\noindent
      \int_compare:nNnTF { \seq_count:N \l__rr_location_pages_seq } = {1}
        { The~corresponding~change~appears~in~the~revised~manuscript~on~ }
        { The~corresponding~changes~appear~in~the~revised~manuscript~on~ }
      \__rr_location_en_pages:.\par
    \fi
    \ifShowChinese
      \par\noindent
      相应修改位于修订稿
      \__rr_location_cn_pages:。\par
    \fi
  }
\NewDocumentCommand{\RenderAutoLocation}{+m}{%
  \__rr_render_auto_location:n {#1}%
}
\ExplSyntaxOff

% -----------------------------------------------------------------------------
% 一条审稿意见的友好字段环境
%
% \begin{ReviewItem}{稳定编号}
%   \CommentEN{完整英文意见}
%   \CommentCN{完整中文意见}
%   \ResponseEN{英文固定首句之后的回复正文}
%   \ResponseCN{中文固定首句之后的回复正文}
%   \RevisedEN{可选：英文修订原文}
%   \RevisedCN{可选：中文对照}
%   \Evidence{可选：图、表或公式}
%   \Location{\At{8}{210--225}\At{8}{244--260}\At{10}{301}\At{12}{365}}
%   特殊位置无法用页码和行号表达时，改用 \LocationEN 和 \LocationCN 手写。
% \end{ReviewItem}
% -----------------------------------------------------------------------------
\NewDocumentEnvironment{ReviewItem}{m +b}{%
  \begingroup
  \long\def\RRCommentEN{xxx}%
  \long\def\RRCommentCN{xxx}%
  \long\def\RRResponseEN{xxx}%
  \long\def\RRResponseCN{xxx}%
  \long\def\RRRevisedEN{}%
  \long\def\RRRevisedCN{}%
  \long\def\RREvidence{}%
  \long\def\RRAutoLocation{}%
  \long\def\RRLocationEN{xxx}%
  \long\def\RRLocationCN{xxx}%

  \long\def\CommentEN##1{\long\def\RRCommentEN{##1}}%
  \long\def\CommentCN##1{\long\def\RRCommentCN{##1}}%
  \long\def\ResponseEN##1{\long\def\RRResponseEN{##1}}%
  \long\def\ResponseCN##1{\long\def\RRResponseCN{##1}}%
  \long\def\RevisedEN##1{\long\def\RRRevisedEN{##1}}%
  \long\def\RevisedCN##1{\long\def\RRRevisedCN{##1}}%
  \long\def\Evidence##1{\long\def\RREvidence{##1}}%
  \long\def\Location##1{\long\def\RRAutoLocation{##1}}%
  \long\def\LocationEN##1{\long\def\RRLocationEN{##1}}%
  \long\def\LocationCN##1{\long\def\RRLocationCN{##1}}%

  #2%

  \RenderCommentBlock{#1}{\RRCommentEN}{\RRCommentCN}%
  \RenderResponseBlock{\RRResponseEN}{\RRResponseCN}%
  \ifdefempty{\RRRevisedEN}{%
    \ifdefempty{\RRRevisedCN}{}{\RenderRevisedText{\RRRevisedEN}{\RRRevisedCN}}%
  }{\RenderRevisedText{\RRRevisedEN}{\RRRevisedCN}}%
  \ifdefempty{\RREvidence}{}{\RREvidence}%
  \ifdefempty{\RRAutoLocation}{%
    \RenderLocation{\RRLocationEN}{\RRLocationCN}%
  }{\RenderAutoLocation{\RRAutoLocation}}%
  \endgroup
}{}

\NewDocumentCommand{\FigurePlaceholder}{O{0.66\linewidth} m m}{%
  \begin{figure}[htbp]
    \centering
    \fbox{\rule{0pt}{30mm}\rule{#1}{0pt}}
    \caption{\ModeInline{#2}{#3}}
  \end{figure}
}

\newcolumntype{Y}{>{\raggedright\arraybackslash}X}

\begin{document}

\MakeLetterHeader

% =============================================================================
%                      用户编辑区 2：逐条审稿意见与回复
% 每个字段均填写一整段或多段内容，不要逐句拆分。
% =============================================================================

\ReviewerSection{1}

\begin{ReviewItem}{R1.1}
  % 审稿意见英文原文（完整粘贴，不要改写）
  \CommentEN{xxx}
  % 审稿意见中文译文（完整填写）
  \CommentCN{xxx}
  % 英文回复正文；固定首句由模板自动添加，只需填写首句之后的内容
  \ResponseEN{%
    Specifically, xxx. Accordingly, we have revised xxx to improve xxx.%
  }
  % 中文回复正文；固定首句由模板自动添加，只需填写首句之后的内容
  \ResponseCN{%
    具体而言，xxx。因此，我们已对 xxx 进行了修改，以进一步完善 xxx。%
  }
  % 可选：修订稿中的英文原文；不需要引用时可删除本行
  \RevisedEN{xxx}
  % 可选：修订原文的中文对照；不需要引用时可删除本行
  \RevisedCN{xxx}
  % 修改位置：每个 \At{页码}{行号} 表示一处；重复 \At 即可添加多个位置
  % 单个位置示例：\Location{\At{8}{210--225}}
  % 多个位置示例（相同页码会自动合并）：
  % \Location{\At{8}{210--225}\At{8}{244--260}\At{10}{301}\At{12}{365}}
  \Location{\At{xxx}{xxx}}
\end{ReviewItem}

\begin{ReviewItem}{R1.2}
  % 审稿意见英文原文（完整粘贴，不要改写）
  \CommentEN{xxx}
  % 审稿意见中文译文（完整填写）
  \CommentCN{xxx}
  % 英文回复正文；固定首句由模板自动添加
  \ResponseEN{%
    We agree that the original presentation did not make this point sufficiently
    clear. We have therefore clarified xxx and added xxx to support this revision.%
  }
  % 中文回复正文；固定首句由模板自动添加
  \ResponseCN{%
    我们同意原稿对此说明得不够清楚。因此，我们已进一步澄清 xxx，并补充 xxx 作为
    支撑。%
  }
  % 可选：回复中需要展示的图、表或公式；不需要时可删除整个 Evidence 字段
  \Evidence{%
    % 替换为真实图片时使用：
    % \includegraphics[width=0.66\linewidth]{figures/your-figure.pdf}
    \FigurePlaceholder
      {xxx}
      {xxx}%
  }
  % 修改位置：重复 \At{页码}{行号或行号范围} 可添加多个位置
  \Location{\At{xxx}{xxx}}
\end{ReviewItem}

\begin{ReviewItem}{R1.3}
  % 审稿意见英文原文（完整粘贴，不要改写）
  \CommentEN{xxx}
  % 审稿意见中文译文（完整填写）
  \CommentCN{xxx}
  % 英文回复正文；固定首句由模板自动添加
  \ResponseEN{%
    The requested comparison has been added and is summarized in
    Table~\ref{tab:comparison}.%
  }
  % 中文回复正文；固定首句由模板自动添加
  \ResponseCN{%
    我们已补充所要求的比较，结果汇总于表~\ref{tab:comparison}。%
  }
  % 可选：定量比较表；将表题、方法名、指标名和数值中的 xxx 替换为实际内容
  \Evidence{%
    \begin{table}[htbp]
      \centering
      \caption{\ModeInline
        {xxx}
        {xxx}}
      \label{tab:comparison}
      \begin{tabularx}{0.86\linewidth}{@{}Y>{\centering\arraybackslash}p{2.6cm}>{\centering\arraybackslash}p{2.6cm}@{}}
        \toprule
        \ModeInline{Method}{方法} & \ModeInline{Metric 1}{指标 1} & \ModeInline{Metric 2}{指标 2} \\
        \midrule
        \ModeInline{Proposed method}{本文方法} & xxx & xxx \\
        \ModeInline{Baseline method}{基线方法} & xxx & xxx \\
        \bottomrule
      \end{tabularx}
    \end{table}%
  }
  % 修改位置：重复 \At{页码}{行号或行号范围} 可添加多个位置
  \Location{\At{xxx}{xxx}}
\end{ReviewItem}

\ReviewerSection{2}

\begin{ReviewItem}{R2.1}
  % 审稿意见英文原文（完整粘贴，不要改写）
  \CommentEN{xxx}
  % 审稿意见中文译文（完整填写）
  \CommentCN{xxx}
  % 英文回复正文；固定首句由模板自动添加
  \ResponseEN{%
    In response, we have revised xxx and added xxx to substantiate this point.%
  }
  % 中文回复正文；固定首句由模板自动添加
  \ResponseCN{%
    作为回应，我们已修改 xxx，并补充 xxx 以提供支撑。%
  }
  % 修改位置：重复 \At{页码}{行号或行号范围} 可添加多个位置
  \Location{\At{xxx}{xxx}}
\end{ReviewItem}

% 添加新意见：复制一个完整的 ReviewItem 环境。
% 添加新审稿人：复制 ReviewerSection，并从新的 ReviewItem 开始。

\end{document}
